\chapter{Аналитическая/Теоретическая часть}

\section{Обзор предметной области}

\subsection{Методы Agentic Deep Graph Reasoning}
Современные подходы к графовому выводу в искусственном интеллекте эволюционируют от пассивных однопроходных моделей к динамическим агентным архитектурам (Agentic Deep Graph Reasoning) \cite{ref1, ref4, ref5, ref6}. В рамках таких архитектур автономный LLM-агент строит и расширяет знания итеративно в процессе решения задачи \cite{ref8, ref11}. Он не только извлекает факты, но и формирует локальный граф рассуждений (reasoning graph), интегрируя его с глобальным, тем самым разрешая противоречия и принимая обоснованные решения. Такие системы сочетают нейронные возможности LLM с четкостью и интерпретируемостью символьных методов (графов знаний), повышая точность и устраняя галлюцинации при ответах \cite{ref2, ref3}. Однако применяемость таких подходов к физическим роботам, где контекст постоянно меняется в физическом времени и пространстве, остается малоизученной.

\subsection{Темпоральные графы знаний (TKG) и алгоритмы работы с ними}
Традиционные статические графы знаний, представляющие данные в виде триплетов (субъект, предикат, объект), не способны моделировать эволюцию среды. Темпоральные графы расширяют эту концепцию добавлением времени в формате четверок (субъект, предикат, объект, время). Для работы с ними существуют различные алгоритмы прогнозирования и вывода (TKG Completion). Они делятся на эмбеддинговые (создание многомерных векторов фактов, чувствительных ко времени, например TComplEx, TRKGE), алгоритмы на основе логических правил и алгоритмы с использованием графовых нейросетей (GNN). Современные методы, такие как EvoReasoner или LLM-DA, показывают значительные успехи в динамической адаптации правил под управлением LLM \cite{ref12, ref22, ref23}. Тем не менее, методы пополнения временных графов часто фокусируются на глобальных исторических базах данных (ICEWS, GDELT) с гранулярностью событий в минуты или сутки и интервальным временем.

\subsection{Инкрементальное обновление и ограничения для мобильной робототехники}
В задаче навигации и управления сервисным роботом данные (например, детектируемые объекты, препятствия, положения человека) поступают непрерывно, с частотой сенсоров (кадры в секунду), и имеют строгую пространственную привязку. Основное ограничение существующих подходов к TKGC и агентному рассуждению заключается в отсутствии средств для быстрого инкрементального (постепенного) обновления графа с сохранением временной связности и разрешением конфликтов в реальном времени, необходимом для мобильных роботов. Если робот обнаруживает, что чашка, которая находилась на столе, теперь перенесена на полку, система должна не перестраивать глобальный граф, а инкрементально добавить событие перемещения объекта, выявить конфликты в текущем плане действий с помощью LLM-агента и адаптировать этот план. Отсутствие подобных средств инкрементального обновления TKG, предназначенных прямо для навигации мобильного робота, обосновывает необходимость решения поставленной технической задачи.

\section{Постановка задачи (Техническое задание)}

Для достижения поставленной цели необходимо разработать программно-архитектурное решение, включающее в себя модель темпорального графа и модуль управления с агентным рассуждением. Исходя из анализа недостатков существующих средств, формируется следующая постановка задачи исследования:

\begin{enumerate}
    \item \textbf{Разработка и формализация модели темпорального графа знаний (TKG):}
    \begin{itemize}
        \item Разработать структуру графа, позволяющую описывать состояние среды, объекты и пространственно-временные отношения между ними, действия робота и логические зависимости.
        \item Обеспечить структуру графа механизмами поддержки версионирования и сохранения истории перемещения объектов.
    \end{itemize}
    \item \textbf{Разработка метода инкрементального обновления графа:}
    \begin{itemize}
        \item Спроектировать алгоритм инкрементального обновления графа данных, поступающих от сенсорных систем робота (например, модулей детекции объектов).
        \item Интегрировать методы разрешения конфликтов в данных (например, когда объект фиксируется в один и тот же момент времени в разных местах из-за погрешности распознавания) без необходимости полного перестроения базы знаний.
    \end{itemize}
    \item \textbf{Адаптация механизмов KGC (Knowledge Graph Completion):}
    \begin{itemize}
        \item Адаптировать методы KGC для прогнозирования недостающих связей и предиктивного анализа будущих состояний динамических объектов (предсказание траекторий или состояний среды).
    \end{itemize}
    \item \textbf{Механизм семантического планирования на основе Agentic Reasoning:}
    \begin{itemize}
        \item Построить LLM-агента, интерпретирующего высокоуровневые естественные языковые команды (например, ``принеси объект X'', ``найди переместившийся объект Y'').
        \item Реализовать алгоритмы динамической корректировки сформированного плана действий в случаях, когда инкрементальное обновление графа выявляет препятствия или изменение местоположения целевых объектов во время выполнения задачи.
    \end{itemize}
    \item \textbf{Программная реализация и экспериментальная оценка:}
    \begin{itemize}
        \item Разработать программный модуль, реализующий спроектированные алгоритмы.
        \item Провести интеграционно-системное тестирование модуля на открытых датасетах или в симулированной среде мобильного робота, оценив успешность выполнения составных задач, скорость инкрементального обновления графа и качество адаптации планов в динамически меняющейся среде по заданным метрикам.
    \end{itemize}
\end{enumerate}
